\chapter{Performance Engineering on One GPU}\label{ch:perf}

A distributed training job is a single-GPU program replicated thousands of times, with
communication added. Every inefficiency inside one device multiplies across the cluster:
a step that leaves half of each GPU idle wastes half the cluster, however clever the
parallelism. This chapter is about making a single device do as much useful work as possible,
and about the measurements you need before you change anything.

The tools come first: honest timers, the Nsight Systems profiler with NVTX annotations, the
PyTorch memory snapshot, and saved-tensor hooks that show exactly what autograd keeps for the
backward pass. With those in hand we take apart a training step: where the time goes, where the
memory goes, what mixed precision and \code{torch.compile} change, and how activation
checkpointing trades FLOPs for bytes. The chapter ends with the largest single optimisation for
long-context training, a FlashAttention-2 kernel in Triton with \emph{both} the forward and the
backward pass, wired into autograd. It builds directly on the derivation in
\cref{sec:attention:flash} and on the forward-only kernel of \cref{sec:triton:flash}.

\begin{figure}[t]
\centering
\resizebox{\linewidth}{!}{%
\begin{tikzpicture}[
  st/.style={draw, rounded corners=3pt, minimum height=1.05cm, text width=2.35cm, align=center,
    font=\small, line width=0.6pt},
  num/.style={circle, fill=sysblue, text=white, font=\bfseries\scriptsize, inner sep=1.6pt},
  arr/.style={-{Stealth[length=2mm]}, line width=0.7pt, draw=black!60}]
\node[st, fill=fdata!15, draw=keygreen] (a) at (0,0) {Workbench \&\\honest timing};
\node[st, fill=fdata!15, draw=keygreen] (b) at (3.0,0) {Nsight Systems\\\& NVTX};
\node[st, fill=fcompute!15, draw=sysblue] (c) at (6.0,0) {Where time goes;\\mixed precision};
\node[st, fill=fmemory!18, draw=napkin] (d) at (9.0,0) {Memory snapshots\\\& residuals};
\node[st, fill=fmemory!18, draw=napkin] (e) at (9.0,-1.75) {Checkpointing,\\offload, compile};
\node[st, fill=fcomm!14, draw=labpurple] (f) at (6.0,-1.75) {Attention at\\long context};
\node[st, fill=fcomm!14, draw=labpurple] (g) at (3.0,-1.75) {FlashAttention-2\\fwd + bwd in Triton};
\node[st, fill=black!6, draw=black!60] (h) at (0,-1.75) {End-to-end\\step playbook};
\foreach \n/\k in {a/1,b/2,c/3,d/4,e/5,f/6,g/7,h/8} \node[num] at (\n.north west) {\k};
\draw[arr] (a)--(b); \draw[arr] (b)--(c); \draw[arr] (c)--(d); \draw[arr] (d)--(e);
\draw[arr] (e)--(f); \draw[arr] (f)--(g); \draw[arr] (g)--(h);
\end{tikzpicture}}
\caption{Roadmap of the chapter. Stages 1--2 build the measurement tools, 3--5 apply them to
compute and memory, 6--7 replace attention with a kernel of our own, and 8 combines everything into
a checklist for a fast training step.}
\label{fig:perf:roadmap}
\end{figure}

% =====================================================================
\section{A performance workbench}\label{sec:perf:workbench}

Performance work is experimental science, and the first deliverable is a laboratory in which
experiments are cheap, repeatable and recorded. Concretely, that means four things.

\paragraph{A pinned, recorded environment.} Driver, CUDA runtime, PyTorch, Triton and NCCL
versions all change kernel selection and therefore timings; a PyTorch upgrade can change which
attention backend runs, and a Triton upgrade can change the language itself (recent Triton
releases removed the block-pointer API in favour of tensor descriptors, see
\cref{sec:perf:fa2:blockptr}). Use a lock file (\code{uv.lock} or \code{requirements.txt} with exact
versions) and record the environment next to every result. \code{code/projects/a2/env\_check.py}
prints versions, GPU models, SM counts, the interconnect topology (\code{nvidia-smi topo -m}),
whether \code{nsys}/\code{ncu} are installed, and a one-second matmul sanity check against the
published peak:

\begin{lstlisting}[style=shell]
$ python projects/a2/env_check.py --json > results/env.json
$ python projects/a2/env_check.py
torch: 2.x.y+cu12z   cuda: 12.z   triton: 3.y   nccl: 2.y.z
gpu[0]: {'name': 'NVIDIA H100 80GB HBM3', 'mem_GiB': 79.1, 'sms': 132, 'cc': '9.0'}
matmul_check: {'bf16_TFLOPs': ..., 'fraction_of_peak': ...}
nsys: /usr/local/bin/nsys   ncu: /usr/local/bin/ncu
\end{lstlisting}

\paragraph{Benchmarks as programs with flags.} Every variation you will want to compare
(model size, context length, dtype, compiled or not, which attention, which checkpointing
policy) should be a command-line flag of one script, not a copy of it. Lab~13's
\code{bench.py} follows this pattern with four sub-commands (\code{model}, \code{attention},
\code{saved}, \code{ckpt}).

\paragraph{Sweeps that log failures.} Out-of-memory errors are results, not exceptions to swallow.
\code{projects/a2/sweep.py} runs a command over a JSON grid, records one JSON line per run
(including \code{"status": "oom"}), stores the git revision, and resumes after an interruption.

\paragraph{Tables generated by code.} Every number in a writeup should come from a file produced by
a script. \code{projects/a2/tables.py} pivots a sweep's JSONL into Markdown or LaTeX, which your
report then \code{\textbackslash input}s. Copying numbers by hand is how mismatched units and stale
results get into papers.

\begin{lstlisting}[style=shell,caption={A sweep from grid to LaTeX table.},label={lst:perf:sweep}]
$ cat grid.json
{"cfg": ["gpt-small", "gpt-medium"], "T": [256, 512, 1024], "dtype": ["fp32", "bf16"]}
$ python projects/a2/sweep.py --grid grid.json --out runs.jsonl -- \
    python labs/lab13_perf/bench.py model --cfg {cfg} --T {T} --dtype {dtype} --jsonl /dev/stdout
$ python projects/a2/tables.py runs.jsonl --rows cfg T --cols dtype --value ms_mean --latex \
    > report/step_time.tex
\end{lstlisting}

% =====================================================================
\section{Timing GPU code honestly}\label{sec:perf:timing}

\Cref{sec:pytorch:async} showed that kernel launches are asynchronous: \code{torch.matmul}
returns as soon as the kernel is \emph{queued}. A timer around it measures launch overhead, not
work. \Cref{lst:perf:timefn} is the timing loop used throughout the chapter (it lives in
\code{llmsys/perf.py}); it extends \cref{lst:pytorch:bench} by returning the full distribution.

\begin{lstlisting}[caption={Warm up, synchronise, repeat, keep every sample.},label={lst:perf:timefn}]
def time_fn(fn, warmup=5, iters=10, sync=synchronize):
    for _ in range(warmup):            # absorb one-off costs
        fn()
    sync()                             # drain everything queued so far
    times = []
    for _ in range(iters):
        t0 = timeit.default_timer()    # highest-resolution monotonic clock
        fn()
        sync()                         # wait for *this* call's kernels
        times.append(timeit.default_timer() - t0)
    return dict(mean=statistics.fmean(times), std=statistics.pstdev(times),
                median=statistics.median(times), min=min(times), times=times)
\end{lstlisting}

\begin{figure}[t]
\centering
\resizebox{\linewidth}{!}{%
\begin{tikzpicture}[x=0.9cm, y=0.55cm, font=\scriptsize]
\node[anchor=east] at (-0.2,2) {CPU (Python)};
\node[anchor=east] at (-0.2,0.6) {GPU stream};
\foreach \x/\w in {0/0.35,0.45/0.35,0.9/0.35,1.35/0.35} \fill[sysblue!45] (\x,1.7) rectangle ++(\w,0.6);
\node[anchor=south] at (0.85,2.35) {launch $k_1..k_4$};
\fill[warnred!35] (1.8,1.7) rectangle ++(0.25,0.6); \node[anchor=south] at (1.93,2.35) {\code{t1}};
\draw[warnred, thick] (1.93,1.6)--(1.93,-0.2); \node[warnred, anchor=north] at (1.93,-0.25) {wrong stop};
\foreach \x/\w/\c in {0.4/1.3/fcompute, 1.75/1.6/fcompute, 3.4/0.9/fdata, 4.35/1.4/fcompute}
  \fill[\c!55] (\x,0.3) rectangle ++(\w,0.6);
\node at (1.05,0.6) {$k_1$}; \node at (2.55,0.6) {$k_2$}; \node at (3.85,0.6) {$k_3$}; \node at (5.05,0.6) {$k_4$};
\fill[keygreen!50] (2.05,1.7) rectangle (5.75,2.3); \node at (3.9,2.0) {\code{synchronize()} blocks};
\draw[keygreen!60!black, thick] (5.8,2.5)--(5.8,-0.2); \node[keygreen!50!black, anchor=north] at (5.8,-0.25) {correct stop};
% warmup
\begin{scope}[xshift=7.3cm]
\node[anchor=west,font=\scriptsize\bfseries] at (0,2.7) {Per-iteration time};
\foreach \i/\h in {1/2.6,2/1.6,3/0.95,4/0.82,5/0.8,6/0.81,7/0.79,8/0.8,9/0.82,10/0.8}
  \fill[\ifnum\i<4 warnred!45\else sysblue!45\fi] (0.45*\i,0) rectangle ++(0.33,\h);
\draw[->] (0.3,0)--(5.0,0) node[right] {iter};
\node[warnred, font=\tiny, align=center] at (1.0,3.0) {warm-up: compile,\\autotune, allocator};
\end{scope}
\end{tikzpicture}}
\caption{Left: the CPU races ahead of the GPU. A timer stopped right after the launches measures
almost nothing; synchronising first makes the CPU wait for the queue to drain. Right
(schematic): the first iterations are slower and must be excluded from the statistics.}
\label{fig:perf:timeline}
\end{figure}

\paragraph{What warm-up absorbs.} The first call to a model on a fresh process pays for things that
never recur: CUDA context creation and module loading (hundreds of milliseconds), cuBLAS/cuDNN
heuristic selection or autotuning for each new shape, Triton JIT compilation and autotuning,
\code{torch.compile} tracing and code generation (seconds to minutes), growth of the caching
allocator's pool (each new block is a \code{cudaMalloc}), and GPU clock ramp-up from idle. If the
first few samples differ from the rest, which they nearly always do, either discard them or
report them separately as ``time to first step'', a number that matters for short jobs and for
serving cold starts.

\paragraph{Variance and what to report.} Report a central value \emph{and} a spread: the mean
with standard deviation if the samples look Gaussian, otherwise the median with percentiles.
Large spread usually has a cause: another process on the GPU, thermal or power throttling
(\code{nvidia-smi -q -d PERFORMANCE} shows the throttle reasons), host-side Python work that is
sometimes on the critical path, or data-dependent work (e.g.\ a variable number of tokens).

\paragraph{CUDA events and \code{do\_bench}.} Wall-clock timing includes host overhead between
synchronisations. CUDA events (\code{torch.cuda.Event(enable\_timing=True)}) are timestamps written
by the GPU into its own stream, so the elapsed time between two events measures device time only
(\code{perf.cuda\_event\_time}). Triton's \code{triton.testing.do\_bench} goes further: it estimates
how many repetitions fit in a time budget, clears the L2 cache between repetitions by writing a
large buffer (so a kernel cannot benefit from data left in cache by its previous run), and returns
a robust statistic. \code{perf.bench} uses it on GPUs and falls back to \code{time\_fn} elsewhere.
Use events or \code{do\_bench} for individual kernels and wall-clock timing for whole training
steps, where host overhead is part of what you want to measure.

\begin{napkinmath}[What should a forward pass cost?]
Before timing anything, predict the answer. GPT-2 medium has $3.53\times10^{8}$ parameters, of
which $3.02\times10^8$ are outside the embeddings. A batch of $B=4$ sequences of $T=512$ tokens is
2048 tokens; with $\approx2N$ FLOPs per token plus the attention term (\code{flops.flops\_per\_token\_forward})
the forward pass needs about $1.55\times10^{12}$ FLOPs. At the H100's 989\,TFLOP/s bf16 peak that is
1.6\,ms; a realistic 40--60\% of peak gives 2.6--3.9\,ms. In fp32 without TF32 the matmuls run on
the CUDA cores at roughly 67\,TFLOP/s, so expect about 15 times longer. A measurement far outside
these bounds means either the measurement or your understanding is wrong, and both are worth
knowing.
\end{napkinmath}

% =====================================================================
\section{Profiling with Nsight Systems and NVTX}\label{sec:perf:nsys}

A timer says \emph{how long}; a profiler says \emph{where}. Python profilers such as
\code{cProfile} see only the host side: they attribute time to the function that \emph{launched}
a kernel, which is the wrong place. NVIDIA Nsight Systems (\code{nsys}) instruments the CUDA
driver through CUPTI and records, on one timeline, the host threads, every CUDA API call, every
kernel and memory copy on the device, NCCL communication, OS runtime events, and user-defined
\emph{NVTX ranges}. Nsight Compute (\code{ncu}) is its microscope: it replays a single kernel many
times to read hardware counters (achieved occupancy, memory throughput per level, tensor-pipe
utilisation, warp stall reasons) and places the kernel on a roofline.

\subsection{Collecting a profile}

\begin{lstlisting}[style=shell,caption={Profiling only the measured iterations.},label={lst:perf:nsys}]
nsys profile -o step_bf16 \
  --trace=cuda,nvtx,osrt,cublas,cudnn \
  --pytorch=functions-trace,autograd-shapes-nvtx \
  --capture-range=nvtx --nvtx-capture=measure --capture-range-end=stop \
  python labs/lab13_perf/bench.py model --cfg gpt-medium --T 512 --dtype bf16 --nvtx
nsys stats --report cuda_gpu_kern_sum step_bf16.nsys-rep     # kernel summary as text
\end{lstlisting}

The flags are worth understanding rather than copying. \code{--trace} chooses which APIs to
record; every additional source adds overhead. \code{--pytorch=functions-trace} makes recent
\code{nsys} versions emit an NVTX range per PyTorch operator and \code{autograd-shapes-nvtx} labels
autograd nodes with tensor shapes; both are convenient, and both slow the run. The capture-range
options start recording only when an NVTX range named \code{measure} opens, so warm-up iterations
(and their compile and autotune kernels) never pollute the profile. Backtrace options
(\code{--cudabacktrace}, \code{--python-backtrace}) tell you which line launched each kernel at a
large overhead: enable them for a targeted question, not by default. The \code{.nsys-rep} file
opens in the Nsight Systems GUI on any laptop, so you can profile on a remote GPU node and inspect
the result locally.

\subsection{Annotating your code}

NVTX ranges are cheap host-side push/pop markers. They appear as labelled bars above the CUDA API
row, and the GUI can filter every statistic by range. \code{llmsys.perf} provides a context manager
that emits both an NVTX range (for Nsight) and a \code{record\_function} scope (for
\code{torch.profiler}), plus a helper that annotates module calls automatically with hooks:

\begin{lstlisting}[caption={NVTX annotation by hand and by hooks (from \texttt{llmsys/perf.py}).},label={lst:perf:nvtx}]
@contextlib.contextmanager
def nvtx_range(name):
    if torch.cuda.is_available():
        torch.cuda.nvtx.range_push(name)
    try:
        with torch.profiler.record_function(name):
            yield
    finally:
        if torch.cuda.is_available():
            torch.cuda.nvtx.range_pop()

def annotate_modules(model, types=()):
    nvtx, handles = torch.cuda.nvtx, []
    for name, mod in model.named_modules():
        if mod is model or (types and not isinstance(mod, tuple(types))):
            continue
        handles += [mod.register_forward_pre_hook(
                        lambda m, a, l=name: nvtx.range_push("fwd:" + l)),
                    mod.register_forward_hook(lambda m, a, o: nvtx.range_pop()),
                    mod.register_full_backward_pre_hook(
                        lambda m, g, l=name: nvtx.range_push("bwd:" + l)),
                    mod.register_full_backward_hook(lambda m, gi, go: nvtx.range_pop())]
    return handles

# inside an attention implementation:
with nvtx_range("attn/scores"):  s = q @ k.transpose(-1, -2) * scale
with nvtx_range("attn/softmax"): p = torch.softmax(s, dim=-1)
with nvtx_range("attn/pv"):      o = p @ v
\end{lstlisting}

To annotate code you do not own, monkey-patch the function before building the model
(\code{some\_module.attention = annotated\_attention}); Python resolves module-level names at call
time. One caveat: \code{torch.compile} fuses across your annotations and may move work out of a
range, so profile compiled code at the granularity of whole compiled regions.

\begin{figure}[t]
\centering
\resizebox{\linewidth}{!}{%
\begin{tikzpicture}[x=0.75cm, y=0.5cm, font=\scriptsize]
\foreach \y/\l in {4/NVTX, 3/CUDA API (host), 2/GPU kernels, 1/memory ops, 0/NCCL}
  \node[anchor=east] at (-0.2,\y+0.3) {\l};
\fill[keygreen!35] (0,4) rectangle (9.2,4.6); \node at (2.0,4.3) {fwd};
\fill[keygreen!55] (9.4,4) rectangle (20.8,4.6); \node at (14.0,4.3) {bwd};
\fill[keygreen!20] (21.0,4) rectangle (23.6,4.6); \node at (22.3,4.3) {opt};
\foreach \x in {0,0.4,...,9.0} \fill[sysblue!50] (\x,3) rectangle ++(0.25,0.6);
\foreach \x in {9.4,9.8,...,20.6} \fill[sysblue!50] (\x,3) rectangle ++(0.25,0.6);
\foreach \x in {21,21.2,...,23.4} \fill[sysblue!50] (\x,3) rectangle ++(0.12,0.6);
% kernels: big gemms and small elementwise
\foreach \x/\w/\c in {0.3/1.4/fcompute,1.8/0.3/fmemory,2.2/0.4/fmemory,2.7/1.4/fcompute,4.2/0.3/fmemory,4.6/1.6/fcompute,6.3/0.3/fmemory,6.7/0.3/fmemory,7.1/1.3/fcompute,8.5/0.5/fmemory}
  \fill[\c!60] (\x,2) rectangle ++(\w,0.6);
\foreach \x/\w/\c in {9.6/1.9/fcompute,11.6/0.4/fmemory,12.1/2.0/fcompute,14.2/0.4/fmemory,14.7/2.6/fcompute,17.4/0.3/fmemory,17.8/1.8/fcompute,19.7/0.9/fcompute}
  \fill[\c!60] (\x,2) rectangle ++(\w,0.6);
\foreach \x in {21.1,21.35,...,23.5} \fill[fmemory!60] (\x,2) rectangle ++(0.15,0.6);
\fill[fcomm!50] (12.0,0) rectangle (16.5,0.6); \node at (14.25,0.3) {all-reduce (overlapped)};
\fill[fdata!50] (0,1) rectangle (0.25,1.6);
\draw[warnred, thick, decorate, decoration={brace, amplitude=3pt}] (21.0,1.85)--(23.6,1.85);
\node[warnred, anchor=north] at (22.3,1.0) {launch-bound gaps};
\end{tikzpicture}}
\caption{What a training step looks like in Nsight Systems (schematic). Host-side launches (CUDA
API row) run ahead of device execution. Large blue GEMM kernels dominate forward and backward;
short orange memory-bound kernels (norms, activations, softmax, elementwise) fill the gaps. In an
unfused optimizer step, hundreds of tiny kernels leave the GPU idle between launches: the
signature of a launch-bound region.}
\label{fig:perf:nsys}
\end{figure}

\subsection{Reading the timeline}

Four questions cover most investigations. \textbf{Is the GPU ever idle?} White space in the
kernel row while the host row is busy means the host cannot launch fast enough (launch-bound:
too many small kernels, Python overhead, a synchronisation such as \code{.item()}); white space in
both rows means the host is waiting for something else (data loading, I/O, a collective).
\textbf{Which kernels dominate?} The \emph{CUDA GPU Kernel Summary} (\code{nsys stats} or the GUI's
Stats view) ranks kernels by total time; filter it by NVTX range to split forward from backward
from optimizer. GEMM kernels have recognisable names (\code{sm90\_xmma\_gemm\_\dots},
\code{cutlass\_\dots}, \code{ampere\_\dots\_gemm}), and Triton kernels carry their Python
function name. \textbf{Does the host match the device?} Clicking a CUDA API call highlights the
kernel it launched, which shows the queue depth: how far ahead the CPU is running. \textbf{Is
communication overlapped?} NCCL kernels concurrent with compute kernels on another stream are
free; NCCL kernels alone on the timeline are not.

\begin{systemslens}[FLOPs are not time]
In a transformer forward pass more than 95\% of FLOPs are in GEMMs, but a profile typically shows
non-GEMM kernels taking a much larger share of time, often 20--40\% for small models in fp32 and more
in bf16, where the GEMMs got faster but the elementwise kernels did not. The reason is arithmetic
intensity (\cref{sec:pytorch:roofline}): a softmax over a $T\times T$ score matrix does a handful of
FLOPs per element and moves every element through HBM at least twice, whereas the matmul that
produced it does $2\dhead$ FLOPs per element. With $\dhead=64$ that is a FLOP ratio of roughly
$50\times$ in the matmul's favour but a \emph{time} ratio that can be close to $1\times$, because both
are paying for the same bytes. Fusion (\cref{sec:perf:residuals}) and FlashAttention
(\cref{sec:perf:fa2}) attack exactly this gap.
\end{systemslens}

% =====================================================================
\section{Where the time goes in a training step}\label{sec:perf:anatomy}

\begin{table}[t]
\centering\small
\begin{tabularx}{\linewidth}{@{}lXll@{}}
\toprule
kernel class & examples & FLOPs/byte (bf16) & usually\\
\midrule
GEMM & QKV/out projections, MLP, LM head & $\approx\frac{mnk}{mk+kn+mn}$, hundreds & compute-bound\\
attention core & $\mat{Q}\mat{K}^\top$, softmax, $\mat{P}\mat{V}$ (unfused) & GEMMs high, softmax $\sim1$ & mixed\\
normalisation & LayerNorm, RMSNorm & $\sim1$--3 & memory-bound\\
activation & GELU, SiLU$\cdot$gate & $<1$--2 & memory-bound\\
residual / dropout / casts & add, mask, \code{.to(bf16)} & $<1$ & memory-bound\\
loss & log-softmax over $V$, gather & $\sim1$ & memory-bound, huge\\
optimizer & AdamW over all parameters & $\approx0.5$ & memory-bound\\
\bottomrule
\end{tabularx}
\caption{Kernel classes in a transformer training step, with typical arithmetic intensity
(analytic). Only the GEMMs can approach the tensor-core roofline; everything else is limited by
HBM bandwidth, so its cost is set by the bytes it moves.}
\label{tab:perf:classes}
\end{table}

\Cref{tab:perf:classes} is the mental model to bring to a profile. Three non-obvious entries
deserve comment.

\paragraph{The loss.} The LM head produces $B\cdot T\cdot V$ logits. For Qwen3's vocabulary
($V=151{,}936$) and a modest batch of $4\times2048$ tokens, that is 2.5\,GB in bf16 and 5.0\,GB once
upcast to fp32 for a stable log-softmax, plus the same again for its gradient. The FLOPs are
trivial; the bytes are not. \Cref{sec:perf:playbook} shows how to avoid materialising them.

\paragraph{The optimizer.} An unfused AdamW step reads the parameter, gradient and two moments and
writes back the parameter and both moments: about 28 bytes per fp32 parameter. For 1.3B parameters
that is 36\,GB of HBM traffic, or 11\,ms at the H100's 3.35\,TB/s \emph{if} the kernels saturate
bandwidth, which tiny per-tensor kernels do not. The forward and backward passes of the same
model on $8\times2048$ tokens need about 130\,ms at peak, so the optimizer is a visible slice of
the step, larger the smaller the batch. PyTorch's \code{foreach} implementation (the default)
batches tensors into a few multi-tensor kernels; \code{fused=True} goes further and does the whole
update in one kernel per parameter group.

\paragraph{Backward is about twice forward.} Each linear layer $\mat{Y}=\mat{X}\mat{W}$ needs
$\mat{dX}=\mat{dY}\mat{W}^\top$ and $\mat{dW}=\mat{X}^\top\mat{dY}$ in backward: two GEMMs of the
same size as the one in forward. Elementwise ops are similar or worse, because their backward
reads saved inputs \emph{and} the incoming gradient. Expect backward time of
$1.8$--$2.5\times$ forward; a much larger ratio points at an inefficient backward kernel (classically
attention, or an op with a scatter in its gradient).

\begin{systemslens}[Tile and wave quantisation]
GEMM throughput is not a smooth function of shape. A GEMM kernel covers the output with tiles
(e.g.\ $128\times256$); a dimension that is one more than a multiple of the tile size launches a
whole extra row of mostly empty tiles (\emph{tile quantisation}). The tiles then run in
\emph{waves} of one tile per SM: 133 tiles on a 132-SM H100 take two waves, the second with a
single busy SM (\emph{wave quantisation}). Hence the folklore of padding $\dmodel$, $\dff$,
vocabulary and per-GPU batch $\times$ sequence to multiples of 64 or 128, and of looking at
throughput as a function of size before choosing a model's shape (Exercise~\ref{ex:perf:waves}).
Memory-access alignment adds a third effect: rows whose byte stride is not a multiple of the
DRAM burst and vector-load width force uncoalesced loads.
\end{systemslens}

% =====================================================================
\section{Mixed precision in practice}\label{sec:perf:amp}

\Cref{sec:pretrain:mixed} introduced the recipe: bf16 matmuls under \code{torch.autocast},
fp32 master weights, loss scaling only for fp16. Here we look at what autocast actually does to
a block, why accumulation precision matters more than storage precision, and what changes in
memory.

\paragraph{Autocast is a per-op policy.} Inside \code{torch.autocast(device\_type, dtype)}, each
operator consults a list. Ops that benefit from tensor cores and are numerically safe in 16 bits
(\code{mm}, \code{addmm}, \code{bmm}, \code{linear}, convolutions, SDPA) cast their inputs to the
low-precision dtype. Ops that are numerically sensitive (reductions such as \code{sum} and
\code{mean}, \code{softmax}/\code{log\_softmax}, \code{layer\_norm}, losses, \code{exp},
\code{pow}) run in fp32. Everything else runs in the widest dtype among its inputs. Parameters
stay fp32; autocast casts a weight once per forward pass and caches the copy, so gradients flow
back into the fp32 parameter.

\begin{figure}[t]
\centering
\resizebox{\linewidth}{!}{%
\begin{tikzpicture}[font=\scriptsize,
  op/.style={draw, rounded corners=2pt, minimum height=0.7cm, minimum width=1.55cm, align=center},
  lo/.style={op, fill=fcompute!25, draw=sysblue},
  hi/.style={op, fill=fmemory!30, draw=napkin},
  arr/.style={-{Stealth[length=1.6mm]}, draw=black!60}]
\node[hi] (x) at (0,0) {$\vect x$\\fp32};
\node[hi] (n1) at (2.0,0) {RMSNorm\\fp32 math};
\node[lo] (qkv) at (4.0,0) {QKV linear\\bf16 in/out};
\node[lo] (att) at (6.0,0) {SDPA\\bf16};
\node[lo] (o) at (8.0,0) {out proj\\bf16};
\node[op, fill=black!8] (r1) at (10.0,0) {$+$ residual\\fp32 (promote)};
\node[hi] (n2) at (12.0,0) {RMSNorm\\fp32 math};
\node[lo] (up) at (14.0,0) {gate/up\\bf16};
\node[op, fill=black!8] (act) at (16.0,0) {SiLU$\cdot$\\bf16};
\node[lo] (dn) at (18.0,0) {down\\bf16};
\foreach \a/\b in {x/n1,n1/qkv,qkv/att,att/o,o/r1,r1/n2,n2/up,up/act,act/dn} \draw[arr] (\a)--(\b);
\node[hi] (loss) at (18.0,-1.5) {log-softmax + CE\\fp32};
\node[lo] (head) at (16.0,-1.5) {LM head\\bf16 logits};
\draw[arr] (dn) -- ++(0,-0.75) -| (head); \draw[arr] (head)--(loss);
\node[anchor=west, align=left] at (0,-1.5) {\tikz\fill[fcompute!25,draw=sysblue] (0,0) rectangle (0.3,0.2);\ cast to bf16 (tensor cores)\quad
  \tikz\fill[fmemory!30,draw=napkin] (0,0) rectangle (0.3,0.2);\ forced fp32\quad
  \tikz\fill[black!8,draw=black] (0,0) rectangle (0.3,0.2);\ widest input dtype};
\end{tikzpicture}}
\caption{dtype flow through a pre-norm block under \code{torch.autocast(dtype=torch.bfloat16)}.
Exact lists differ by PyTorch version and backend (check \code{torch.amp.autocast\_mode}); the
pattern is stable: matmuls go to bf16, reductions and losses stay fp32, and the fp32 residual
stream promotes bf16 branch outputs back to fp32.}
\label{fig:perf:amp}
\end{figure}

\paragraph{Accumulate wide, store narrow.} The danger in low precision is not storing a number in
16 bits; it is adding many small numbers into a 16-bit running sum. Once the sum is large, each
addend is smaller than half a unit in the last place of the sum and is rounded away. A
measurement in this environment: adding $10^{-3}$ ten thousand times gives 9.9945 with an fp32
accumulator (the error is just bf16's rounding of $10^{-3}$ itself), 4.0 with an fp16 accumulator,
and 0.5 with a bf16 accumulator. bf16 has fp32's \emph{range} but only 8 bits of \emph{precision},
so it stagnates sooner than fp16. This is why tensor cores multiply 16-bit inputs but accumulate in
fp32, why every reduction in a kernel (the row sums $\ell$ and the output accumulator in
FlashAttention, the variance in RMSNorm) should be kept in fp32 registers, and why
\code{torch.sum} on a bf16 tensor accumulates internally in fp32.

\paragraph{fp16 versus bf16 for normalisation.} With fp16, the mean of squares in an RMSNorm can
overflow (squares of activations above $\sim256$ exceed fp16's 65{,}504) and small variances can
underflow; this is why autocast runs norms in fp32. bf16 removes the range problem but not the
precision one, so the statistic should still be computed in fp32, which is what
\code{llmsys.layers.RMSNorm} does explicitly.

\begin{pitfall}[Mixed precision can \emph{increase} memory]
Autocast keeps fp32 parameters \emph{and} caches bf16 copies of the weights during the forward
pass, and the bf16 activations saved for backward are only half of what they would be in fp32 if the
saved tensors are actually bf16. Some are not: the fp32 residual stream and the inputs of
fp32-forced ops stay fp32. For small models, where weights rather than activations dominate memory,
peak memory under autocast can be equal to or slightly higher than in pure fp32. The large memory
wins come from \emph{storing} parameters in bf16 (with fp32 master copies held only by the
optimizer, or sharded as in \cref{ch:sharding}) and from bf16 activations at long contexts.
Measure it with \cref{sec:perf:memory}'s tools instead of assuming it.
\end{pitfall}

\begin{lstlisting}[caption={Precision switches worth knowing.},label={lst:perf:ampknobs}]
torch.backends.cuda.matmul.allow_tf32 = True     # fp32 GEMMs on tensor cores (10-bit mantissa)
torch.set_float32_matmul_precision("high")        # same, newer spelling
with torch.autocast("cuda", dtype=torch.bfloat16):
    loss = loss_fn(model(x), y)                   # matmuls bf16, reductions fp32
loss.backward()                                   # grads land in fp32 .grad
# fp16 needs dynamic loss scaling:
scaler = torch.amp.GradScaler("cuda")
scaler.scale(loss).backward(); scaler.step(opt); scaler.update()
\end{lstlisting}

